\documentclass[12pt, a4paper]{article}
\usepackage{silence}
\WarningFilter{latex}{You have requested package}
\RequirePackage{../../lectures}
\usepackage[left=3cm,right=3cm,top=3cm,bottom=3cm]{geometry}
\usepackage{mathtools}
\usepackage{float}


\title{Конспект лекций по алгебре и геометрии}
\author{Смирнов Дмитрий Станиславович sdsik@yandex.ru}
\date{}

\begin{document}
    \maketitle
    \tableofcontents
    \section{Комплексные числа}

    \begin{definition}
        Комплексное число - выражение вида \( a + bi \), где \( a, b \in \mathbb{R} \), а \( i \)- особый символ, причем \( i^2 = -1 \).\\
        В записи выше \(  a\) называется вещественной (действительной) частью комплексного числа, а \( b \) - мнимой. \( a = \mathrm{Re}\,z, b = \mathrm{Im}\, z \)
    \end{definition}
    \subsection{Операции над комплексными числами}
    Пусть \( z_{1} = a_{1} + b_{1}i \), \( z_{2} = a_{2} + b_{2}i \) - комплексные числа.
    \subsubsection{Равенство комплексных чисел}
    По определению 
    \[
        z_{1} = z_{2} \iff
    \begin{cases}
        a_{1} = a_{2} \\
        b_{1} = b_{2}
    \end{cases}
    \] 
    \subsubsection{Сложение комплексных чисел}
    \[
        z_{1} + z_{2} \coloneqq (a_{1} + a_{2}) + (b_{1} + b_{2})i
    \] 

    \begin{exmp}
        \( 3 - 2i \) и \( 4 + 3i \)

        +: \( 7 + i \)

        -: \( -1 - 5i \)
    \end{exmp}

    \subsubsection{Умножение комплексных чисел}
    Умножение происходит как у многочленов по переменной \( i \):
    \[
        z_{1} \cdot z_{2} = (a_{1} + b_{1}i) \cdot (a_{2} + b_{2}i)= (a_{1}a_{2} - b_{1}b_{2}) + (a_{1}b_{2} + b_{1}a_{2})i
    \] 
    
    
    \subsubsection{Сопряжение комплексных чисел}
    \begin{definition}
        Сопряженным для комплексного числа \( z = a + bi \) будет называться комплексное число \( \overline{z} = a - bi \)
    \end{definition}

    \begin{exmp}
        \item 1. \( \overline{z} = 3z \)
            \begin{align*}
                \overline{z} &= 3z \\
                a - bi &= 3a + 3bi \\
                2a + 4bi &= 0 \\
                a + 2bi &= 0 \\
                 \Updownarrow \\
            \begin{cases*}
                a = 0 \\
                b = 0 
            \end{cases*} 
            \end{align*} 
            То есть, это верно только при \( z = 0 \).
        \item 2. \( z^2 = \overline{z} \)
            \begin{align*}
                z^2 & = \overline{z} \\
                a^2 + 2abi - b^2 & = a - bi \\
                \Updownarrow \\
                \begin{cases} 
                    a^2 - b^2  = a \\
                    2ab = - b \\
                \end{cases}
            \end{align*}
            Разбираем второе уравнение: \( 2ab = -b \iff 2ab + b = 0 \iff b(2a + 1) = 0 \iff b = 0 \) или \( a = - \frac{1}{2} \)

            Подставляем \( b = 0 \) в первое уравнение: 
            \begin{align*}
                a^2 & = a \\
                    & \Updownarrow \\
                \left[ \begin{array}{l}
                        a = 1\\
                        a = 0
                \end{array}\right.
            \end{align*}
            Подставляем \( a = -\frac{1}{2} \) в первое уравнение:
            \begin{align*}
                \frac{1}{4} - b^2 &= \frac{1}{2} \\
                \frac{3}{4} = b^2 &\implies b = \pm \frac{\sqrt{3} }{2} 
            \end{align*}
            Собираем все вместе: \( z = 0; 1; -\frac{1}{2} + \frac{\sqrt{3} }{2}; -\frac{1}{2} - \frac{\sqrt{3} }{2} \)
    \end{exmp}

    \subsubsection{Деление комплексных чисел}
    Для деления комплексных чисел нужно домножить знаменатель на сопряженное: \[
        \frac{z_{1}}{z_{2}} = \frac{a_{1} + b_{1}i}{a_{2} + b_{2}i} = \frac{(a_{1} + b_{1}i)(a_{2} - b_{2}i)}{(a_{2} + b_{2}i)(a_{2} - b_{2}i)} = \frac{(a_{1} + b_{1}i)(a_{2} - b_{2}i)}{a_{2}^2 + b_{2}^2} 
    \] 
    Числитель - комплексное, знаменатель - заведомо действительное, следовательно, можно разделить числитель покомпонентно.

    \begin{exmp}
    \item 1. \( \frac{3}{2 - i} \)
        \[
            \frac{3}{2 - i} = \frac{3(2 + i)}{5} = \frac{6}{5} + \frac{3}{5}i
        \] 
        
    \item 2. \( \frac{-3 + 2i}{1 + i} \)
        \[
            \frac{-3 + 2i}{1 + i} = \frac{(-3 + 2i)(1 - i)}{2} = \frac{-3 + 3i + 2i + 2}{2} = -\frac{1}{2} + \frac{5}{2}i
        \] 
    \end{exmp}

    \subsubsection{Возведение в степень}
    \[
        z^{n} = \underbrace{z \cdot z \cdot \ldots \cdot z }_{ n \text{ раз}}
    \] 
    \[
        (1 + i)^{20} = ((1 + i)^2)^{10} = (12 + 2i - 1)^{10} = (2i)^{10} = -1024 
    \] 
    \subsection{Квадратные корни из комплексного числа в алгебраической форме}
    \begin{definition}
        Пусть \( z \) - комплексное число, тогда квадратным корнем из \( z \) называют число, квадрат которого равен \( z \)
    \end{definition}
    Вычисление удобно производить через метод неопределенных коэффициентов: \[
        \sqrt{z} = A + Bi \implies z = (A + Bi)^2
    \] 
    \begin{exmp}
        \begin{align*}
            \sqrt{8 + 6i}  &= A + Bi \\
            8 + 6i &= A^2 - B^2 + 2ABi
        \end{align*}
        Приравниваем действительные и мнимые части: \[
            \begin{cases}
                8 = A^2 - B^2 \\
                6 = 2AB
            \end{cases} \implies A = \frac{6}{2B} = \frac{3}{B}
        \] 
        Подставляем в первое уравнение и получаем биквадратное уравнение:
        \begin{align*}
            8B^2 - 9 + B^{4} &= 0 \\
            B = \pm 1 \implies A &= \pm 3
        \end{align*}
        Ответ: \( 3 + i \); \( -3 - i \)
    \end{exmp}
    \subsection{Квадратные уравнения с комплексными коэффициентами и корнями}
    \begin{align*}
        &az^2 + bz + c = 0 \\
        &z_{1, 2} = \frac{-b \pm \sqrt{b^2 - 4ac} }{2a}
    \end{align*}
    Поскольку для комплексных чисел арифметический квадратный корень не определен в этом случае под \( \pm \sqrt{b^2 - 4ac}  \) подразумевается прибавление квадратных корней из дискриминанта. Для ненулевого дискриминанта таких чисел два и они противоположны по знаку. Для извлечения квадратного корня из дискриминанта применяют вычисления, аналогичные приведенным выше.
    \begin{exmp}
        \[
            z^2 + (-1 + i)z + 2 - 2i = 0
        \]. \( D = b^2 - 4ac = (-1 + i)^2 - 4(2 - 2i) = 1 - 2i - 1 -8 + 8i = -8 + 6i \)
        \[
            \begin{cases}
                A^2 - B^2 = -8 \\
                2AB = 6
            \end{cases} \implies A = \frac{3}{B} \implies \frac{9}{B^2} - B^2 = -8
        \] 
        \begin{align*}
            &9 - B^{4}  = 8B^2 \\
            &B^{4} - 8B^2 - 9 = 0\\  
        \end{align*}
        \begin{align*}
            B = \pm 3 \\
            A = \pm 1
        \end{align*} \( \implies \sqrt{D} = \pm(1 + 3i) \)

        \begin{align*}
            z_{1} = \frac{1 - i + 1 + 3i}{2} = 1 + i \\
            z_{2} = \frac{1 - i -1 - 3i}{2} = -2i
        \end{align*}
        Ответ: \( 1 + i \); \( -2i \)
    \end{exmp}
    \subsection{Модуль и аргумент комплексного числа}

    \begin{figure}[htbp]
        \incfig{complex_geometry}
        \caption{Комплексное число как вектор.}
    \end{figure}
    Длина \( |z| \) -- модуль комплексного числа.

    \( \mathrm{Arg} z  = \varphi\)  -- аргумент комплексного числа.

    \( \mathrm{Arg} z = \arg z + 2 \pi k, \, k \in \mathbb{Z}, \, 0 \le \arg z < 2 \pi \)

    По теореме Пифагора: \[
        |z| = \sqrt{x^2 + y^2} 
    \] 

    Также по определению косинуса и синуса:
    \[
        \begin{cases}
            \cos  \varphi = \frac{x}{\sqrt{x^2 + y^2} } \\
            \sin  \varphi = \frac{y}{\sqrt{x^2 + y^2} }
        \end{cases}
    \] 

    Теперь можно переписать определение комплексного числа в следующей форме : \[
        z = x + iy = \sqrt{x^2 + y^2} \left( \frac{x}{\sqrt{x^2+y^2} } + \frac{y}{\sqrt{x^2 + y^2} }i \right)  = |z| ( \cos \varphi + i \sin \varphi)
    \] 
    Самое правое выражение называют тригонометрической формой комплексного числа. \( |z| \) иногда обозначают \( r \).

    Из формул Эйлера для синуса и косинуса: \[
        \begin{cases}
            \cos \varphi = \frac{e^{i\varphi} + e^{-i\varphi}  }{2} \\
            \sin  \varphi = \frac{e^{i\varphi} - e^{-i\varphi}  }{2i}
        \end{cases}
    \] следует еще и показательная формула комплексного числа: \[
        z = re^{i\varphi}. 
    \] 

    \begin{exmp}
        \(  \)
        \begin{figure}[htbp]
            \incfig{trig_form_example}
            \caption{Пример.}
        \end{figure}

        \begin{enumerate}
            \item \( 2 = 2(\cos 0 + i \sin 0) \)
            \item \( -6 = 6(\cos  \pi  + i \sin  \pi ) \)
            \item \( i = 1 (\cos  \frac{\pi}{2} + i \sin \frac{\pi}{2})  \)
            \item \( 1 - i = \sqrt{2} (\cos \frac{7 \pi }{4} + i \sin \frac{7 \pi }{4}) \)
            \item \( i + i\sqrt{3}  = 2 (\cos \frac{\pi}{3} + i \sin \frac{\pi}{3}) \)
        \end{enumerate}
    \end{exmp}

    \subsection{Умножение и возведение в степень в тригонометрической форме}
    Пусть \( z_{1} = r_{1}(\cos \varphi_{1} + i \sin \varphi_{2}) \), \( z_{2} = r_{2}(\cos \varphi_{2} + i \sin \varphi_{2}) \), тогда
    \begin{align*}
        z_{1}\cdot z_{2} &= r_{1}r_{2}(\cos \varphi_{1} + i \sin \varphi_{1}) (\cos \varphi_{1} + i \sin \varphi_{2}) = \\
                         &=r_{1}r_{2}(\cos \varphi_{1}\cos \varphi_{2} + i \sin\varphi_{1}\cos \varphi_{2} + i \sin \varphi_{2}\cos \varphi_{1} - \sin \varphi_{2}\cos \varphi_{2}) = \\
                         & = r_{1}r_{2}(\cos (\varphi_{1} + \varphi_{2}) + i \sin(\varphi_{1} + \varphi_{2}))
    \end{align*}
    Аналогично выводится формула для деления: \[
        \frac{z_{1}}{z_{2}} = \frac{r_{1}}{r_{2}}(\cos (\varphi_{1} - \varphi_{2}) + i \sin (\varphi_{1} - \varphi_{2}))
    \] 
    Теперь попробуем умножить \( z \) само на себя:
    \[
        z^2 = r^2(\cos  2\varphi + i \sin 2\varphi)
    \] 
    Из этого следует (желающие могут доказать по индукции) \[
        z^{n} = r^{n}(\cos n\varphi + i \sin n\varphi)  
    \] 
    \subsection{Формулы Муавра}
    В принципе, то, что мы получили чуть выше и есть формула Муавра. Но в исходных текстах она выглядит так: \[
        (\cos \varphi + i \sin \varphi)^{n} = \cos n\varphi + i \sin \varphi 
    \] 
    \begin{exmp}
        Вычислим формулы для \( \cos 5x \) и \( \sin 5x \) через \( \sin x \) и \( \cos  x \)
        \begin{lem}
            \( (a + b)^{5} = a^{5} + 5a^{4}b + 10a^{3}b^2 + 10a^2b^3 + 5ab^{4} + b^{5}  \)
            
        \end{lem}
        \begin{proof}
            Гуглить: бином Ньютона, треугольник Паскаля
        \end{proof}

        Из формулы Муавра: \( (\underline{\cos 5 \varphi}+ i\underline{\underline{ \sin 5 \varphi}}) = (\cos  \varphi + i \sin  \varphi)^{5}  \). Раскроем правую скобку при помощи леммы. 
        \begin{align*}
            (\cos  \varphi + i \sin  \varphi)^{5} = \cos ^{5}\varphi &+ 5 \cos  ^{4}(i \sin \varphi)  + 10 \cos  ^3 \varphi(i \sin  \varphi)^2 + 10 \cos ^2\varphi(i \sin  \varphi)^3  + \\
            &+ 5 \cos \varphi (i \sin \varphi)^{4} + (i \sin \varphi)^{5} = \\
            & =  \underline{\cos ^{5}\varphi - 10 \cos ^3 \varphi \sin ^2\varphi + 5 \cos \varphi \sin ^{4} \varphi} + \\
            & + i (\underline{\underline{5 \cos ^{4} \varphi \sin \varphi - 10 \cos ^2\varphi \sin ^3 \varphi + \sin ^{5} \varphi}}  )
        \end{align*}
        Осталось только приравнять действительные (одно подчеркивание) и мнимые (два подчеркивания) части равенства:
        \begin{align*}
            \cos 5 \varphi = \cos ^{5}\varphi - 10 \cos ^3 \varphi \sin ^2\varphi + 5 \cos \varphi \sin ^{4} \varphi \\
            \sin 5\varphi = 5 \cos ^{4} \varphi \sin \varphi - 10 \cos ^2 \varphi \sin ^3 \varphi + \sin  ^{5}\varphi  
        \end{align*}
        
    \end{exmp}

    \begin{exmp}
        Упростим сумму \( \sin x + \sin  2x + \ldots + \sin nx \).
        \begin{align*}
            z &= \cos  \varphi + i \sin \varphi \\
            z^{n} &= \cos n\varphi + i \sin n\varphi 
        \end{align*}
        Несложно заметить, что \[
            \sin x + \sin 2x + \ldots  + \sin  nx = \mathrm{Im}(z + z^2 + \ldots + z^{n} )
        \] 
        По формуле суммы геометрической прогрессии:\[
            z + z^2 + \ldots  + z^{n} = \frac{z(1 - z^{n} )}{1 - z}.
        \] Подставив вместо \( z \) тригонометрическую форму, и проведя ряд несложных преобразований, получаем \[
            \sin x + \sin 2x + \ldots + \sin nx = \frac{\sin \frac{nx}{2} \cdot \sin \frac{n+1}{2}x}{\sin \frac{x}{2}}
        \] 
        
    \end{exmp}
    \subsection{Корни из комплексного числа в тригонометрической форме}
    \begin{definition}
        Пусть \( n \in \mathbb{N} \). Корнем степени \( n \) из комплексного числа \( z \) называется такое число \( x \), что \( x^{n} = z  \).
    \end{definition}
    \begin{align*}
        z &= r(\cos \varphi + i \sin \varphi) \\
        x' &= p(\cos \beta + i \sin \beta) \\
        (x')^{n} &= p^{n}(\cos n\beta + i \sin n\beta) = r (\cos \varphi + i \sin \varphi) = z
    \end{align*}
    Отсюда получаем \[
        \begin{cases}
            p^{n} = r \\
            \cos n\beta = \cos \varphi\\
            \sin n\beta = \sin \varphi
        \end{cases}
    \] 
    Так как равны одновременно и синусы и косинусы, то углы отличаются на \( 2 \pi \):
    \begin{align*}
        n\beta &= \varphi + 2 \pi k, \qquad k \in \mathbb{Z}\\
        \beta &= \frac{\varphi + 2 \pi k}{n}\\
        p &= \sqrt[n]{r} \\
        x' &= \sqrt[n]{r}\left( \cos \frac{\varphi + 2 \pi k}{n} + i \sin \frac{\varphi + 2 \pi k}{n} \right) , \qquad k = 0, 1, 2, \ldots , n -1 
    \end{align*}
    Корни степени \( n \) из числа \( z \) равномерно расположены на окружности радиуса \( \sqrt[n]{r}  \).
    \begin{exmp}
        \( \sqrt[3]{8}  \)
        \begin{align*}
            \sqrt[3]{8} &= \sqrt[3]{8(\cos 0 + i \sin 0)} = 2 \left( \cos \frac{2 \pi k}{3} + i \sin \frac{2 \pi k}{3} \right), \qquad k = 0, 1, 2 \\
            k &= 0 \quad 2\left( \cos 0 + i \sin 0 \right) = 2 \\
            k &= 1 \quad 2\left( \cos \frac{2 \pi }{3} + i \sin \frac{2 \pi }{3}\right) = 2\left(- \frac{1}{2} + i \frac{\sqrt{3} }{2}\right)  \\
            k &= 2 \quad 2\left( \cos \frac{4 \pi }{3} + i \sin \frac{4 \pi }{3}\right) = 2\left(- \frac{1}{2} - i \frac{\sqrt{3} }{2}\right)  
        \end{align*}
        
    \end{exmp}

    \begin{exmp}
        \( \sqrt[4]{i}  \)
        \begin{align*}
            \sqrt[4]{i} &= \sqrt[4]{\cos \frac{\pi}{2} + i \sin \frac{\pi}{2}}  = \cos \frac{\frac{\pi}{2} + 2 \pi k}{4} + i \sin \frac{\frac{\pi}{2} + 2 \pi k}{4}, \qquad k = 0, 1, 2, 3\\
            k &= 0 \quad \cos \frac{\pi}{8} + i \sin \frac{\pi}{8}\\
            k &= 1 \quad \cos \frac{5\pi}{8} + i \sin \frac{5\pi}{8}\\
            k &= 2 \quad \cos \frac{9\pi}{8} + i \sin \frac{9\pi}{8}\\
            k &= 3 \quad \cos \frac{13\pi}{8} + i \sin \frac{13\pi}{8}\\
        \end{align*}
    \end{exmp}
    \subsection{Изображение множеств на плоскости}
    \begin{enumerate}
        \item \( |z - z_{1}| = |z - z_{2}| \)

            Другими словами, это \( \sqrt{(x - x_{1})^2 + (y - y_{1})^2}  = \sqrt{(x - x_{2})^2 + (y - y_{2})^2}  \), то есть, множество точек, равноудаленных от точек \( z_{1} \) и \( z_{2} \). Из курса школьной математики известно, что это серединный перпендикуляр
            \begin{figure}[ht!]
                \incfig{complex_drawing1}
            \end{figure}
        \item \( |z - z_{1}| \le |z - z_{2}| \)
            
            Аналогичными рассуждениями получаем, что это множество точек, которое ближе к \( z_{1} \), чем к \( z_{2} \), т.е. полуплоскость слева от серединного перпендикуляра
            \begin{figure}[ht]
                \incfig{complex_drawing2}
            \end{figure}
        \item \( |z_{1} \cdot z - z_{2}| \le R \)
            
            Вынесем \( z_{1} \) за скобки: \( |z_{1}| \left| z - \frac{z_{1}}{z_{1}} \right| \le  R \), поделим на \( |z_{1}| \): \( \left| z - \frac{z_{2}}{z_{1}} \right| \le \frac{R}{|z_{1}|}  \). Слева очевидно комплексное число, а справа очевидно действительное. Сделаем замену: \( |z - z_{0}| \le R_{0} \). Теперь проделаем процедуру из предыдущих пунктов: \[
                \sqrt{(x - x_{0})^2 + (y - y_{0})^2}  \le  R_{0}
            \] или же, \[
            (x - x_{0})^2 + (y - y_{0})^2 \le R_{0}^2
            \]   
            А это буквально круг с центром в \( (x_{0}, y_{0}) \) и радиусом \( R_{0} \).
    \end{enumerate}
    \section{Многочлены и действия над ними}
    \begin{definition}
        Многочлен - функция вида \[
            f(x) = a_{n}x^{n} + a_{n-1}x^{n-1} + \ldots  + a_{1}x + a_{0} = \sum_{i=0}^{n} a_i x^{i}.
        \]При этом, \( a_{n} \) называется старшим членом, \( a_{0} \) -- свободным членом, \( n = \deg f \) -- степенью многочлена. Многочлен равен нулю \( \iff \) все его коэффициенты равны нулю. Если в записи многочлена нет какой-либо степени неизвестного, то коэффициент при этой степени равен нулю.
    \end{definition}
    На множестве многочленов определены следующие действия: 
    \begin{enumerate}
        \item Сложение
        \item Умножение 
        \item Деление с остатком

            \[
                \frac{f(x)}{g(x)} \iff f(x) = g(x) \cdot q(x) + r(x), \quad \deg r < \deg g
            \] 
    \end{enumerate}
    \subsection{Корни многочлена}
    \begin{definition}
        Корнем многочлена \( f(x) \) называется число \( a \): \( f(a) = 0 \)
    \end{definition}
    \begin{theorem}[Безу]
        \( \forall f(x), a \in \mathbb{R} \): \[
            f(x) = (x-a) \cdot q(x) + r, \qquad r = f(a)
        \] 
    \end{theorem}
    \begin{proof}
        Доказывается подстановкой \( a \) в \( f \): пусть \( x = a \implies f(a) = r \).
    \end{proof}
    \begin{cor}
        \( a \) -- корень \( f(x) \iff f(x) \, \vdots \, (x - a) \) без остатка.
    \end{cor}
    \subsection{Рациональные корни многочлена с целыми коэффициентами}
    \begin{theorem}
        Пусть \( f(x) \) -- многочлен с целыми коэффициентами. Если \( r = \frac{m}{n} \), несократима дробь, -- корень \( f \), то \(a_{0} \, \vdots \, m \) и \( a_{n} \, \vdots  \,n \).
    \end{theorem}
    \begin{proof}
        Пусть \( a_{0}, a_{1}, \ldots , a_{n} \) -- целые коэффициенты \( F \). \[
            F(x) = a_{n}x^{n} + a_{n-1}x^{n-1} + \ldots  + a_{1}x + a_{0}  .
        \]. Пусть \( r = \frac{p}{q} \) -- несократимая и \( F(\frac{p}{q}) = 0 \), тогда \( a_{0} \, \vdots \, p, \, a_{n} \, \vdots \, q \). Действительно, подставим \( r \) в \( F \):
        \[
            a_{n}\left( \frac{p}{q} \right) ^{n} + \ldots + a_{1} \frac{p}{q} + a_{0} = 0.
        \] Домножим на знаменатель: \[
            a_{n}p^{n} + a_{n-1}p^{n-1}q + \ldots  + a_{1}pq^{n-1} + a_{0}q^{n} = 0.
        \]Теперь внимательно смотрим на то, что делится на \( p \): \[
        \underbrace{a_{n}p^{n} + a_{n-1}p^{n-1}q + \ldots  + a_{1}pq^{n-1}}_{\vdots \, p} + a_{0}q^{n} = \underbrace{0}_{\vdots\, p}.
        \]Ага, то есть, и \( a_{0}q^{n}  \) делится на \( p \). Но т.к. \( \frac{p}{q} \)  -- несократимая, то \( p \) и \( q \) -- взаимно простые, а значит именно \( a_{0} \) делится на \( p \). Абсолютно точно также доказывается то, что \( a_{n} \) делится на \( q \).
        
    \end{proof}
    \begin{rem}
        Теорема дает только необходимое условие существования рационального корня и оно не является достаточным.
    \end{rem}
    \begin{rem}
        Среди делителей нужно брать не только положительные, но и отрицательные целые числа.
    \end{rem}
    \begin{rem}
        Начать проверку предполагаемых корней удобнее всего с небольших целых чисел.
    \end{rem}
    \begin{exmp}
        \( x^3 + 3x^2 + 3x + 2 = 0 \).
        По теореме выше кандидатами на корни являются \( \pm 1, \pm 2 \). Проверяем их: \( 1 \) -- нет. \( -1 \) -- нет. \( 2 \) -- нет. \( -2: -8 + 12 - 6 + 2 = 0 \), то есть \( -2 \) -- корень.
    \end{exmp}
    \subsection{Разложение многочленов на множители}
    \begin{definition}
        Разложить многочлен на множители -- представить его в виде произведения двух многочленов степень каждого из которых не меньше единицы.
    \end{definition}
    \begin{exmp}
        \(  \)

        \begin{enumerate}
            \item \( x^{4} - 1 = (x^2 - 1)(x^2 + 1) = (x^2+1)(x - 1)(x+1) \)
            \item \( x^{4} - x^3 - x + 1 = x^3 (x-1) - (x-1) = (x-1)(x^3 -1) = (x-1)^2(x^2 + x + 1)  \)
            \item \( x^{4} + 4 = x^{4} + 4x^2 + 4 - 4x^2 = (x^2 + 2)^2 - 4x^2 = (x^2 + 2 - 2x)(x^2 + 2 + 2x)   \)
            \item \( x^{8} + x + 1 = x^{8} - x^2 + x^2 + x + 1 = x^2(x^{6} - 1) + x^2 + x + 1 = x^2(x^3  - 1)(x^3 + 1) + x^2 + x + 1 =x^2(x-1)(x^2 + x + 1)(x^3 +1) + x^2 + x + 1 = (x^2 + x + 1)(x^2(x^3  + 1)(x-1) + 1) \)
            \item \( x^{4} + x^2+ 1 = (x^2 + 1)^2 - x^2 = (x^2 - x + 1)(x^2 + x + 1) \)
        \end{enumerate}
    \end{exmp}
    \subsection{Основная теорема алгебры}
    \begin{theorem}
        Каждый многочлен, отличный от константы, имеет хотя бы один комплексный корень.
    \end{theorem}
    \begin{cor}
        Для любого \( n \in \mathbb{N} \) всякий многочлен степени \( n \) имеет ровно \( n \) комплексных корней с учетом кратности, т.е. корень кратности \( k \) считается как \( k \) корней. При этом число \( a \) -- корень кратности \( k \), если многочлен делится на \( (x-a)^{k}  \) и не делится на \( (x-a)^{k+1}  \).
        
    \end{cor}

    \begin{rem}
        Доказать следствие основной теоремы алгебры можно применив следствие из теоремы Безу несколько раз. Сначала ищем первый корень \( x_{1} \), затем делим многочлена на \( (x - x_{1}) \), и для полученного многочлена снова найдем корень. Так повторяем до тех пор, пока не получим многочлен первой степени.
        
    \end{rem}
    \subsection{Корни многочлена с действительными коэффициентами}
    \begin{lem}
        Корни многочлена с действительными коэффициентами, не являющиеся вещественными, разбиваются на пары сопряженных чисел.
        
    \end{lem}
    \begin{proof}
        Рассмотрим многочлен \( a_{n}x^{n} + a_{n-1}x^{n-1} + \ldots + a_{1}x + a_{0}   \), где \( a_{0}, a_{1}, \ldots , a_{n} \in \mathbb{R} \). Если \( z \) -- корень, то \( a_{n}z^{n} + a_{n-1}z^{n-1} + \ldots a_{1}z + a_{0} = 0 \). Подставим теперь в многочлен \( \overline{z} \): \[
            a_{n}\overline{z}^{n} + a_{n-1}\overline{z}^{n - 1} + \ldots + a_{1}\overline{z} + a_{0} = \overline{a_{n} z^{n} + a_{n-1}z^{n-1} + \ldots + a_{1}z + a_{0}  } = \overline{0} = 0.
        \] То есть, если \( z \) -- корень, то и \( \overline{z} \) -- тоже корень.   
    \end{proof}
    \begin{cor}[Разложение на множители степени не выше двух]
        Многочлен с вещественными коэффициентами можно разложить на множители степени не выше второй над \( \mathbb{R} \):
        \begin{align*}
            a_{n}x^{n} + a_{n-1}x^{n-1} + \ldots &+ a_{1}x + a_{0} = \\
                                                                  &= a_{n}(x - x_{1})\ldots (x - x_{k})(x^2 + b_{1}x + c_{1})\ldots (x^2 + b_{s}x + c_{s}),
        \end{align*} где \( b_{i} < 4c_{i}, \, 1 \le i \le s \)
    \end{cor}
    \subsection{Метод Ньютона}
    \begin{definition}
        Метод Ньютона, также известный как метод касательных -- итерационный численный метод нахождения корня заданной функции. Поиск решения осуществляется путем построения последовательных приближений.\[
            x_{n+1} = x_{n} - \frac{f(x_{n})}{f'(x_{n})}.
        \] Уравнение выше следует из уравнения касательной к функции \( f \) в точке \( (x_{0}, y_{0}) \): \( y - y_{0} = y'(x_{0})(x - x_{0}) \).
    \end{definition}
    При этом \( x_{n+1} \) дадет более точное приближение корня, чем \( x_{n} \). И сходимость довольно быстрая.
    \begin{figure}[H]
        \incfig{newthon_method}
    \end{figure}
    \begin{rem}
        \(  \)
    \begin{enumerate}
        \item Я не понял, что он рассказывал в качестве метода Ньютона, так что взял описание из википедии.
        \item Начальное приближение \( (x_{1}, y_{1}) \) должно быть достаточно близким. При плохом приближении метод может и не сойтись.  
    \end{enumerate}
    \end{rem}
    Метод может использоваться и для приближенного вычисления комплексных корней многочлена, правда для комплексного случая нужно дополнительно уточнить область сходимости, а она порой имеет очень своеобразную форму.

    \subsection{Теорема Виета для многочленов степени выше второй}
    Рассмотрим многочлен
    \begin{align*}
        (x - x_{1})(x - x_{2})(x - x_{3}) & = (x^2 - (x_{1} + x_{2})x + x_{1}x_{2})(x-x_{3}) = \\
                                          & = x^3 - (x_{1} + x_{2} + x_{3})x^2 + (x_{1}x_{2} + x_{2}x_{3} + x_{1}x_{3})x - x_{1}x_{2}x_{3} = \\
                                          & = x^3  + bx^2 + cx + d \implies \begin{cases}
                                            x_{1} + x_{2} + x_{3} = -b \\
                                            x_{1}x_{2} + x_{1}x_{3} + x_{2}x_{3} = c \\
                                            x_{1}x_{2}x_{3} = -d
                                          \end{cases}
    \end{align*}
    Очевидно, это можно обобщить и на высшие степени:\[
        x^{n} + a_{n-1}x^{n-1} + \ldots + a_{1}x + a_{0}; \qquad x_{1},x_{2},\ldots , x_{n} \text{ -- корни}  
    \]
    \[
        \begin{cases}
            a_{n-1} = - (x_{1} + x_{2} + \ldots  + x_{n}) \\
            a_{n-2} = x_{1}x_{2} + x_{1}x_{3} + \ldots + x_{1}x_{n} + x_{2}x_{3} + \ldots  + x_{n-1}x_{n}\\
            a_{n-3} = -(x_{1}x_{2}x_{3} + x_{1}x_{2}x_{4} + \ldots + x_{n-2}x_{n-1}x_{n})\\
            \ldots \\
            a_{1} = (-1)^{n-1} (x_{1}x_{2}\ldots x_{n-1} + x_{1}x_{2}\ldots x_{n-2}x_{n} + \ldots  + x_{2}\ldots x_{n-2}x_{n-1}x_{n})\\
            a_{0} = (-1)^{n}(x_{1}x_{2}\ldots x_{n}) 
        \end{cases}
    \] 

    \begin{exmp}
        \( x^3  + 2x^2 - 3x + 4 = 0 \)
        \[
            \begin{cases}
                x_{1} + x_{2} + x_{3} = -2 \\
                x_{1}x_{2} + x_{1}x_{3} + x_{2}x_{3} = -3 \\
                x_{1}x_{2}x_{3} = -4
            \end{cases}
        \] 
        Норешая такую систему в лоб, придем к тому же самому уравнению.
        
    \end{exmp}
    \begin{rem}
        Доказательство, как и для многочлена второй степени, осуществляется рассмотрением равенства, полученного разложением многочлена по корням.
    \end{rem}
    \begin{rem}
        Если старший коэффициент многочлена не равен 1, то для применения теоремы Виета нужно поделить многочлен на этот коэффициент.
    \end{rem}

    \section{Дробно-рациональные функции}
    \begin{definition}
        Дробно рациональная функция (ДРФ) -- отношение \( f = \frac{A(x)}{B(x)}, \quad B(x) \neq 0 \). При этом \( D(f) = \mathbb{R} \backslash Z \), где \( Z \) -- множество всех корней \( B \).
    \end{definition}
    \subsection{Операции над ДРФ}
    \begin{enumerate}
        \item  Равенство. \[
            \frac{A(x)}{B(x)} = \frac{C(x)}{D(x)} \iff A(x) \cdot D(x) = B(x) \cdot C(x)
        \] 
        \item Сложение. \[
                \frac{A(x)}{B(x)} + \frac{C(x)}{D(x)} = \frac{A(x)\cdot D(x) + B(x) \cdot C(x)}{B(x) \cdot D(x)}
        \] 
        \item Умножение. \[
            \frac{A(x)}{B(x)} \cdot \frac{C(x)}{D(x)} = \frac{A(x) \cdot  C(x)}{B(x) \cdot D(x)}
        \]  
    \end{enumerate}
    \subsection{Расширение и сокращение ДРФ}
    \[
        \frac{A(x)}{B(x)} \iff \frac{A(x) \cdot C(x)}{B(x) \cdot C(x)}
    \] 

    \( \Leftarrow  \) Сокращение.

    \( \Rightarrow \) Расширение.
    \begin{exmp}
        \[
            \frac{2x}{x^2 - 1} + \frac{3x-2}{x + 1} = \frac{2x + (3x - 2)(x-1)}{x^2 - 1} = \ldots 
        \] 
        
    \end{exmp}
    \subsection{Првильные и простейшие ДРФ}
    \begin{definition}
        Правильная ДРФ - такая ДРФ, что степень числителя строго меньше степени знаменателя: \( \deg A < \deg B \)
    \end{definition}
    \begin{definition}
        Простейшая ДРФ - функция следующего вида: \( \frac{A(x)}{(B(x))^{n} } \), где \( \deg A < \deg B \), \( n \in N \), \( B(x) \) -- неприводимый над заданным множеством.
    \end{definition}
    \begin{exmp}
        \[
            \underbrace{\frac{x + 2}{(x+1)^3 }}_{\text{правильная, но не простейшая}} = \frac{x + 1 + 1}{(x + 1)^3 } = \underbrace{\frac{1}{(x+1)^2} + \frac{1}{(x+1)^3 }}_{\text{простейшие}}
        \] 
    \end{exmp}
    \begin{lem}
        Любую правильную дробь можно разложить в сумму простейших, причем единственным образом.
        Пусть \( \deg A < \deg B \). \( \frac{A(x)}{B(x)}\), \( B(x) = \prod_{i=1}^{m} (B_{i}(x))^{m}    \). Тогда \[
            \frac{A(x)}{B(x)} = \sum_{i=1}^{m} \sum_{j=1}^{n_{i}} \frac{A_{ij}(x)}{(B_{i}(x))^{j}},
        \] где \( \frac{A_{ij}(x)}{(B_{i}(x))^{j}} \) -- простейшая дробь.
    \end{lem}
    \begin{proof}
        Верь мне, брат
    \end{proof}
    \subsection{Метод неопределенных коэффициентов}
    \begin{exmp}
        Я не понял, что препод хотел этим сказать, видимо, о том,как найти остаток от деления на какой-то из корней, попутно вычисляя коэффициента частного.
        Рассмотрим многочлен из теоремы Безу: \[
            A(x) = (x - x_{0})B(x) + A(x_{0})
        \] 
        Пусть теперь \[
            A(x) = \sum_{i=0}^{m} a_{i}x^{i}, \qquad a_i, x_0\text{ -- известны} 
        \] 
        \[
            B(x) = \sum_{i=0}^{m-1} b_{i}x^{i}, \qquad b_{i} \text{ -- неизвестны}  
        \] 
        \[
            A(x_{0}) \text{ -- неизвестно и требуется}
        \] Подставим эти суммы в исходное равенство: \[
            \sum_{i=0}^{m} a_{i}x^{i}(x - x_{0})  \cdot  \sum_{i=0}^{m-1} b_{i}x^{i} + A(x_{0}) = \ldots  
        \] Раскроем скобки и приведем подобные: \[
            \begin{cases}
                b_{m-1} = a_{m}\\
                b_{m-2}=x_{0} \cdot b_{m-1} + a_{m-1}\\
                b_{m-3} =x_{0} \cdot b_{m-2} +a_{m-2}\\
                \ldots \\
                b_{m-1-i} = x_{0} \cdot b_{m-i} + a_{m-i}, \quad i \in [0, m-1]\\
                b_{0}= x_{0}\cdot b_{1} + a_{1}\\
                A(x_{0}) =x_{0}\cdot b_{0} + a_{0}
            \end{cases}
        \]  
        
    \end{exmp}
    \begin{exmp}[Разложить на простейшие над \( \mathbb{R} \)]
        \[
            \frac{1}{x^2 - 5x + 6} = \frac{1}{(x-2)(x-3)} = \frac{A}{x-2} + \frac{B}{x-3} = \frac{A(x-3) + B(x-2)}{(x-2)(x-3)}
        \] И т.к. знаменатели одинаковые, можем приравнять числители: \[
            1 = A(x-3) + B(x-2) = Ax - 3A  + Bx - 2B = x(A + B) - 3A - 2B.
        \] Приравняем коэффициенты при соответствующих степенях \( x \): \[
            \begin{cases}
                A + B = 0 \\
                -3A - 2B = 1
            \end{cases} \implies B = 1, \, A = -1.
        \] Теперь можно записать ответ: \[
            \frac{1}{x^2 - 5x + 6} = \frac{-1}{x-2} + \frac{1}{x-3}
        \] 
        
    \end{exmp}
    \begin{exmp}[Разложить на простейшие]
        \begin{align*}
            \frac{1}{x^3  + x} & = \frac{A}{x} + \frac{Bx + C}{x^2 + 1} =\\
                                                                      &=\frac{A(x^2 + 1) + x(Bx + C)}{x^3  + x} = \frac{Ax^2 + A + Bx^2 + Cx}{x^3 +x} = \\
                                                                      &= \frac{x^2(A + B) + x\cdot C + A}{x^3  + x}.
        \end{align*}
        Снова составляем систему уравнений и решаем ее: \[
            \begin{cases}
                A = 1 \\
                C = 0 \\
                A + B = 0
            \end{cases} \implies \begin{cases}
                A = 1 \\
                B = -1 \\
                C = 0
            \end{cases}
        \] И запишем ответ: \[
            \frac{1}{x^3 + x} = \frac{1}{x} - \frac{x}{x^2 + 1}.
        \]  
        
    \end{exmp}
    \section{Матрицы}
    \subsection{Множество матриц}
    Множество матриц размера \( m \times n \) обозначается как \( M_{m\times n}(\mathbb{R}) \), где \( m \) -- количество строк, \( n \) -- количество столбцов, \( \mathbb{R} \) -- множество, из которого выбираются элементы матрицы.
    \begin{exmp}
        \(  \)
        \begin{enumerate}
            \item \( A = \begin{pmatrix} 1  & 2 \\ 3 & 4 \end{pmatrix}  \)
            \item \( B = \begin{pmatrix} 1 & 2 & 3 \\ 4 & 5 & 6\end{pmatrix}\)
            \item \( C = \begin{pmatrix} 1 & 2 \\ 3 & 4 \\ 5 & 6 \end{pmatrix}   \)
            \item Матрицы размера \( 1\times 1 \) мы будем отождествлять просто с числами: \( M_{1\times 1}(\mathbb{R}) \equiv \mathbb{R} \)
        \end{enumerate}
        
    \end{exmp}
    Если нам не выжны компоненты матрицы, мы будем обозначать ее так: \( A = (a_{ij})_{m \times n} \).
    \subsection{Операции над матрицами}
    Пусть \( A = (a_{ij})_{m \times n} \), \( B = (b_{ij})_{m \times n} \)
    \begin{enumerate}
        \item Равенство матриц. \( A = B \iff \forall i \in [1, m], \, \forall j \in [1, n] \, a_{ij} = b_{ij} \).  То есть, матрицы равны тогда и только тогда, когда равны все их компоненты.
        \item Сложением матриц. \( A + B = C = (c_{ij})_{m\times n} \), где \( c_{ij} = a_{ij} + b_{ij} \). Соответственно -- поэлементная сумма.
        \item Умножение на скаляр. \( \alpha \cdot A = B \), где \( b_{ij} = \alpha \cdot a_{ij} \).
    \end{enumerate}

    \subsection{Произведение матриц}
    Во-первых, нельзя умножать все подряд матрицы. Они должны быть согласованы по размеру: \[
        C = A \cdot B \qquad A = (a_{ij})_{m\times \underline{\underline{n}}}, \qquad B = (b_{ij})_{\underline{\underline{n}}\times p}.
    \] Где матрица \( C \) определяется как \[
        C = (c_{ik})_{m\times p} \qquad c_{ik} = \sum_{j=1}^{n} a_{ij}b_{jk}
    \] 
    \begin{exmp}
        \[
            \underbrace{\begin{pmatrix} 1 & -1 & 2 \\3 & 0 & 1 \end{pmatrix}}_{2\times 3} \cdot \underbrace{\begin{pmatrix} -1& 3 \\ 4 & 5 \\ -2 & 2 \end{pmatrix}}_{3\times 2}  = A
        \] 
        \( A = (a_{ij}) _{2\times 2} \)\\
        \( a_{11} = 1 \cdot (-1) + (-1) \cdot 4 + 2 \cdot (-2) = -9 \)\\
        \( a_{12} = 1 \cdot 3 + (-1) \cdot 5 + 2 \cdot  2 = 2 \)\\
        \( a_{21} = 3 \cdot (-1) + 0 \cdot 4 + 2 \cdot (-1) = -5 \) \\
        \( a_{22} = 3 \cdot 3 + 0 \cdot 5 + 1 \cdot 2 = 11 \) \\
        \( \implies A = \begin{pmatrix} -9 & 2 \\ -5 & 11 \end{pmatrix}  \)
        
    \end{exmp}
    \subsection{Свойства операций}
    \begin{enumerate}
        \item Ассоциативность сложения. \( (A + B) + C = A + (B + C) \)
        \item Коммутативность сложения. \( A + B = B + A \)
        \item Существование нуля. \( A + O = O + A = A \), здесь \( O \) -- нулевая матрица соответсвующего размера.
    \end{enumerate}
    Заметим, что в общем случае \( A \cdot B \neq B \cdot A \). (Здесь \( m = n \)). То есть, например, \( (A + B)^2 = (A + B)(A + B) = A^2 + AB + BA + B^2 \).
    С учетом некоммутативности выполняются следующие свойства: 
    \begin{enumerate}
        \item Ассоциативность умножения. \( (A \cdot B) \cdot C = A \cdot (B \cdot C) \)
        \item Левая дистрибутивность. \( A \cdot (B + C) = AB + AC \)
        \item Правая дистрибутивность. \( (A + B) \cdot C = AC + BC \)
    \end{enumerate}
    \subsection{Транспонирование матриц и его свойства}
    \begin{definition}
        Транспонирование матрицы -- операция вида \[
            A = (a_{ij})_{m\times n} \implies A^{\mathrm{T}} = (a_{ji})_{n\times m}.
        \] Здесь буквально происходит перестановка столбцов и строк местами.
    \end{definition}
    \begin{exmp}
        \[
            \begin{pmatrix} 1 & -1 \\ 2 & -3 \\ 5 & 4\end{pmatrix}^{\mathrm{T}}  = \begin{pmatrix} 1 & 2 & 5 \\ -1 & -3 & 4 \end{pmatrix} 
        \] 
    \end{exmp}
    Свойства: \begin{enumerate}
        \item Инволютивность. \( (A^{\mathrm{T}} )^{\mathrm{T}} = A  \).
        \item Линейность. \( (\alpha \cdot A + \beta \cdot B)^{\mathrm T} = \alpha \cdot A^{\mathrm T} + \beta \cdot B^{\mathrm T} \).
        \item \( (A \cdot B)^{\mathrm T} = B^{\mathrm T}   \cdot  A^{\mathrm T}   \).
    \end{enumerate}
    \begin{proof}
        \(  \)
    \begin{enumerate}
        \item
            Докажем, что \( B = (A^{\mathrm T} ) ^{\mathrm T} \implies B = A  \). \\
            Пусть \( A = (a_{ij}) _{m \times n} \). Тогда, по определению \( A^{\mathrm T} = (a_{ji}) _{n\times m}  \), и снова по определению \( (A^{\mathrm T} )^{\mathrm T} = (a_{ij})_{m\times n} = B \)
        \item 
            Пусть \( L \coloneqq (\alpha A + \beta B)^{\mathrm T}  \), \( R \coloneqq \alpha A ^{\mathrm T} + \beta B^{\mathrm T} \). Преобразуем оба выраженния и покажем, что они равны.

            \( L = (l_{ij}) _{m \times n} \) 
            \begin{align*}
                \alpha A + \beta B &= (\alpha a_{ij} + \beta b_{ij})_{m \times n} \\
                (\alpha A + \beta B)^{\mathrm T} &= (\alpha a_{ji} + \beta b_{ji}) _{n\times m}.
            \end{align*} То есть, \( l_{ij} = \alpha a_{ji} + \beta b_{ji} \). Теперь с правой частью: 
            \begin{align*}
                \alpha A^{\mathrm T} & = (\alpha a_{ji})_{n\times m} \\
                \beta B^{\mathrm T} & = (\beta b_{ji}) _{n \times m} \\
                \alpha A ^{\mathrm T} + \beta B ^{\mathrm T} & = (\alpha a_{ji} + \beta b_{ji}) _{n \times m} = (r_{ij})_{n\times m}
            \end{align*}
            Итого: \( r_{ij} = \alpha a_{ji} + \beta b_{ji} = l_{ij} \) -- буквально определение равенства матриц.
        \item 
            Пусть \( L \coloneqq (A\cdot B)^{\mathrm T}  \), \( R \coloneqq B^{\mathrm T} \cdot A ^{\mathrm T}\), \( A = (a_{ij})_{m\times n} \), \( B = (b_{jk})_{n\times p} \). Положим \[
                C = AB = \left( \underbrace{\sum_{j=1}^{n} a_{ij} b_{jk}}_{c_{ik}} \right) _{m\times p}
            \] 
            По определению транспонирования \[
                l_{ik} = c_{ki} = \sum_{j=1}^{n} a_{kj}\cdot b_{ji}
            \] 
            Теперь с правой частью: \( B^{\mathrm T} = (b_{kj})_{p\times n}  \), \( A^{\mathrm T} = (a_{ji})_{n\times m}  \) \[
                D = B^{\mathrm T}A^{\mathrm T} = \left( \underbrace{\sum_{j=1}^{m} b_{kj} a_{ji}}_{d_{ki}} \right)   
            \] 
            Теперь замечаем, что \[
                d_{ki} = r_{ki} = c_{ik} = l_{ki},
            \] что снова приводит к определению поэлементного равенства.
            
    \end{enumerate}
    \end{proof}

    \subsection{Элементарные преобразования строк матрицы. Метод Гаусса.}
    \begin{definition}
        Элементарные преобразования строк матрицы: 
        \begin{enumerate}
            \item Перестановка местами двух различных строк \( i \) и \( j \): \( P_{ij} \).
            \item Умножение на ненулевое число \( a \) элементов строки \( i \): \( M_{i}(a) \)
            \item Сложение строки \( i \) с другой строкой \( j \), умноженной на ненулевое число \( a \): \( M_{ij}(a) \)
        \end{enumerate}
    \end{definition}

    \[
        A \xrightarrow{M_{ij}(a)}B \implies \underbrace{B_{k}}_{k-\text{я строка матрицы} B} = \begin{cases}
            A_{k}, \quad k \neq  i \\
            A_{i} + a A_{j}, \quad k =i
        \end{cases}
    \] Рассмотрим систему линейных уравнений(СЛУ): \[
        \begin{cases}
            a_{11}x_{1} + a_{12}x_{2} + \ldots  + a_{1n}x_{n} = b_{1} \\
            a_{21}x_{1} + a_{22}x_{2} + \ldots  + a_{2n}x_{n} = b_{2} \\
            \ldots \\
            a_{m1}x_{1} + a_{m2}x_{2} + \ldots  + a_{mn}x_{n} = b_{m}
        \end{cases}
    \] и перепишем ее в матричном виде. \( A = (a_{ij})_{m\times n} \), \( \overline{x} = (x_{j})_{n\times 1} \), \( \overline{b} = (b_{i})_{m\times 1} \). Тогда \[
    A\overline{x} = \overline{b}.
    \] Здесь \( A \) -- матрица коэффициентов СЛУ, а \( (A \mid  \overline{b}) \) -- расширенная матрица коэффициентов СЛУ.

    Для решения СЛУ, нужно привести расширенную матрицу к особому виду при помощи элементарных преобразований: \[
        (A \mid \overline{b}) \sim  \begin{pmatrix} E & C \\ 0 & 0 \end{pmatrix} 
    \] 

    \begin{definition}
        Ведущий элемент строки матрицы -- первый ненулевай элемент этой строки.
    \end{definition}
    \begin{definition}
        Матрица стандартного вида -- либо матрица, состоящая из одной строки, либо матрица, для каждой строки которой, начиная со второй, ведущий элемент строки находится правее, чем ведущий элемент предшествующих строк, а все нулевые строки матрицы располагаются ниже ненулевых.
    \end{definition}
    \begin{exmp}[Привести к стандартному виду матрицу]
        \[
            \begin{pmatrix} 2 & 3 & -1 \\ 3 & -2 & 4 \end{pmatrix} \xrightarrow{M_{21}\left( -\frac{3}{2} \right) } \begin{pmatrix} 2 & 3 & -1 \\ 0 & -\frac{13}{2} & \frac{11}{2} \end{pmatrix}.
        \] 

        
    \end{exmp}


    
\end{document}
